\documentclass[10pt,a4paper]{amsart}
\usepackage[margin=19mm]{geometry}
\usepackage{amsmath,amssymb,mathtools}
\usepackage[scr=rsfs,cal=euler]{mathalfa}
\usepackage{xfrac}
\usepackage[unicode,hidelinks,bookmarks=false]{hyperref}
\newcommand{\ie}{i.e.\ }
\newcommand{\cf}{cf.\ }
\newcommand{\resp}{respectively\ }
\newcommand{\Subsection}{\subsection}
\providecommand{\nobreakdash}{\nobreak\hbox{-}}
\setlength{\emergencystretch}{2em}
\multlinegap0pt
\usepackage{enumitem}
\usepackage{mathrsfs}
\allowdisplaybreaks%
 \newtheorem{theorem}{Theorem}
 \newtheorem{corollary}{Corollary}
 \newtheorem{lemma}{Lemma}% 
\newtheorem{remark}{Remark}
\newtheorem{definition}{Definition}
\newtheorem{example}{Example}
\begin{document}
\title[On a conjecture concerning antipodal labelings for cycles]{On a conjecture concerning antipodal labelings for cycles}
\author{John M. Campbell}
\date{}
\maketitle
\noindent\emph{Source and licence.} On a conjecture concerning antipodal labelings for cycles. Version arXiv:2609.29477v1 (CC BY 4.0). This material is distributed under the Creative Commons Attribution 4.0 International licence, http://creativecommons.org/licenses/by/4.0/, as recorded in the arXiv OAI licence metadata for this exact version and shown on the abstract page https://arxiv.org/abs/2609.29477v1. Full rights notice: licenses/arxiv-2609.29477-CC-BY-4.0.txt.\par\smallskip
\noindent\textit{Editorial scope.} Counted unit: the original \texttt{theorem} environment \texttt{maintheorem} at source lines 186--188 together with its complete proof at lines 190--434; the delivered document keeps source lines 101--435 so that every referenced label and every citation resolves inside it. Native editable source is the paper's own concatenated TeX; the AMS wrapper is editorial formatting only. No publisher class, upstream script or unrelated artwork is redistributed.\par\smallskip\noindent\textit{Reference note.} This article ships no compiled bibliography (biblatex); the entries delivered here are composed from the author's own BibTeX (.bib) fields, with no field invented and no entry dropped.
\begin{equation}\label{antipodalcondition}
 |f(u) - f(v)| \geq \operatorname{diam}(G) - d(u, v)
\end{equation}
 holds for distinct vertices $u$ and $v$ in $G$. For such a labeling $f$, the \emph{span} of $f$ is then defined so that 
\begin{equation}\label{displayspan} 
 \operatorname{sp}(f) = \max\{ f(u) - f(v) : u, v \in V(G) \}. 
\end{equation} 
 Being consistent with the notation and terminology in the work on Juan and Liu \cite{JuanLiu2012}, 
 the \emph{antipodal number} $\operatorname{an}(G)$
 for $G$ is then defined as the minimum span for an antipodal labeling of $G$. 

 As expressed by Juan and Liu \cite{JuanLiu2012}, the concept of an antipodal labeling for a graph was introduced by Chartrand et al.\ 
 \cite{ChartrandErwinZhang2000,ChartrandErwinZhang2002}, who established general bounds for antipodal numbers. 
 Bounds for $\operatorname{an}(C_n)$ were established by Chartrand et al.\ \cite{ChartrandErwinZhang2000}, and evaluations for 
 $\operatorname{an}(C_{n})$ have been proved for each of the cases among
 $n \equiv 1 \bmod {4}$, 
 $n \equiv 2 \bmod {4}$, and 
 $n \equiv 3 \bmod {4}$ \cite{ChartrandErwinZhang2000,JuanLiu2012}.

 Juan and Liu \cite[Theorem 9]{JuanLiu2012} proved that 
\begin{equation}\label{C4kbounds}
 2k^2 - \left\lfloor \frac{k}{2} \right\rfloor \leq \operatorname{an}(C_{4k}) \leq 2k^2 - 1 
\end{equation}
 for each integer $k$ exceeding $1$, and then conjectured \cite[Conjecture 1]{JuanLiu2012} that $\operatorname{an}(C_{4k}) $ is equal to 
 the upper bound in \eqref{C4kbounds} for every positive integer $k$. Juan and Liu verified this conjecture for $k \in \{ 1, 2, \ldots, 5 \}$. 
 The conjectured evaluation for $\operatorname{an}(C_{n})$ for the $n \equiv 0 \bmod {4}$ case has remained the last unsolved case, out 
 of the congruence classes modulo $4$. 
 This gives weight to our 
 proof of Juan and Liu's conjecture in Section \ref{secmain} below. 
 
\section{Proof of Juan and Liu's conjecture}\label{secmain}
 From the upper bound in \eqref{C4kbounds} proved by Juan and Liu \cite[Theorem 9]{JuanLiu2012}, it remains to prove that the 
 reverse inequality
\begin{equation}\label{displayreverse} 
 \operatorname{an}(C_{4k}) \geq 2k^2 - 1
\end{equation}
 holds for positive integers $k$, 
 and this provides the basis for our proof of Theorem \ref{maintheorem} below. 

\begin{remark}\label{correction}
 Kumar and Panigrahi \cite{KumarPanigrahi2026} do not claim to prove Juan and Liu's conjecture, 
 but they give a related minimality criterion in their Theorem~2.4 and apply it in Theorem~2.5, 
 but it can be shown, as follows, that these results are not valid as stated. Indeed, for \(C_{8}=z_{0}z_{1}\cdots z_{7}z_{0}\), the 
 ordering \(z_{0},z_{4},z_{1},z_{5},z_{2},z_{6},z_{3},z_{7}\), with respective labels \(0,0,3,3,6,6,9,9\), satisfies the hypotheses of Theorem~2.4(a) 
 in Kumar and Panigrahi's paper, 
 as well as those of Theorem~2.5 for \(n=4\), \(d=4\), and \(m=m'=1\), although the resulting span is \(9\). By contrast, the cyclic label sequence \(0,5,2,7,0,5,2,7\) is an antipodal labeling of \(C_{8}\) with span \(7\), and Juan and Liu \cite[Theorem~9]{JuanLiu2012} give \(\operatorname{an}(C_{8})=7\).
\end{remark}

 Let $V(C_{4k}) = \{ v_{0}, v_{1}, \ldots, v_{4k-1} \}$, 
 with the understanding that subscripts are given by residue classes modulo $4k$. 
          Observe that $\operatorname{diam}(C_{4k}) = 2k$. We require 
 the following specialization of 
 a result from Juan and Liu \cite[Proposition 1]{JuanLiu2012}, 
 referring the interested reader to Juan and Liu's proof of a generalized version of the following result. 

\begin{lemma}\label{lemmadistance}
 For any three vertices $u$, $v$, and $w$ in $C_{4k}$, the relation 
 $ d(u, v) + d(v, w) + d(w, u) \leq 4 k $
 holds (cf.\ \cite[Proposition 1]{JuanLiu2012}). 
\end{lemma}

 By taking a given antipodal labeling of a graph, and by then subtracting the minimum label from each label, this does not change the 
 antipodal condition in \eqref{antipodalcondition} or the value of the right-hand side of \eqref{displayspan}. So, for an arbitrary antipodal 
 labeling $f$ of $C_{4k}$, we may assume without loss of generality that $0$ is among the labels. We proceed to fix an ordering on 
 $V(C_{4k})$, by writing 
\begin{equation}\label{VC4k} 
 V(C_{4k}) = \{ x_0, x_1, \ldots, x_{4k-1} \}, 
\end{equation}
 with 
\begin{equation}\label{orderfx} 
 0 = f(x_0) \leq f(x_1) \leq \cdots \leq f(x_{4k-1}). 
\end{equation}
 Borrowing notation from Juan and Liu \cite{JuanLiu2012}, we define $d_{i} = d(x_{i}, x_{i+1})$ and 
\begin{equation}\label{definefi} 
 f_{i} = f(x_{i+1}) - f(x_{i})
\end{equation}
 for $i \in \{ 0, 1, \ldots, 4k-2 \}$. We also require the following specialization of a lemma from Juan and Liu \cite[Lemma 2]{JuanLiu2012}. 
 Again, we refer the interested reader to Juan and Liu's paper for a derivation of a more general version of 
 the following lemma. 
 
\begin{lemma}\label{boundfsum}
 The relation $f_{i} + f_{i+1} \geq k$ holds for all indices 
 $i \in \{ 0, 1, \ldots, 4k - 3 \}$ (cf.\ \cite[Lemma 2]{JuanLiu2012}). 
\end{lemma} 


\begin{theorem}\label{maintheorem}
 The relation $\operatorname{an}(C_{4k}) = 2k^2 - 1$ holds for every positive integer $k$. 
\end{theorem}

\begin{proof}
 To prove the reverse inequality shown in \eqref{displayreverse}, we begin with the base case for $ k = 1$. For an arbitrary antipodal 
 labeling of $C_{4}$, since $\operatorname{diam}(C_4) = 2$, we find that any two adjacent vertices $u$ and $v$ necessarily satisfy $| f(u) - 
 f(v) | \geq \operatorname{diam}(C_4) - d(u, v) = 1$, 
 so that $\operatorname{sp}(f) \geq 1$. 
 Similarly, by assigning the label sequence $0$, $1$, $0$, $1$
 in a cyclic order to the vertices of $C_{4}$, we find that adjacent vertices differ by $1$, whereas opposite vertices
 may have equal labels, giving us an antipodal labeling of span $1$. 
 We thus find that $\operatorname{an}(C_{4k}) = 2k^2 - 1$ holds for $k = 1$. 
 We proceed to disregard this base case and to let $k \geq 2$. 

 Our construction, relying on the ordering in \eqref{orderfx}, gives us that 
\begin{equation}\label{evaluatespf} 
 \operatorname{sp}(f) = f(x_{4k-1}) = \sum_{i=0}^{4k-2} f_{i}, 
\end{equation} 
 recalling the definition displayed in \eqref{definefi}. 
 Moreover, since $f$ is antipodal, we have that 
\begin{equation}\label{2kminusdi} 
 f_{i} \geq 2 k - d_i 
\end{equation}
 for all $i \in \{ 0, 1, \ldots, 4k - 2 \}$. 

 We proceed to define the nonnegative integer
\begin{equation}\label{defineepsilon} 
 \varepsilon_{i} = f_{i} + f_{i+1} - k 
\end{equation}
 for $i \in \{ 0, 1, \ldots, 4k-3 \}$, giving the excess over the lower bound given in Lemma \ref{boundfsum}. From the antipodal condition in 
 \eqref{antipodalcondition}, together with 
 the ordering imposed in \eqref{orderfx}, 
 we deduce that 
\begin{equation}\label{withoutabs}
 f(x_{i + 2}) - f(x_{i}) \geq 2 k - d(x_i, x_{i+2}). 
\end{equation}
 From the definition of $\varepsilon_i$ in \eqref{defineepsilon} together
 with \eqref{withoutabs}, we find that 
 $ k + \varepsilon_i \geq 2 k - d(x_i, x_{i+2}) $
 for all indices $i \in \{ 0, 1, \ldots, 4 k - 3 \}$, 
 i.e., so that 
\begin{equation}\label{negepsilon} 
 d(x_{i}, x_{i+2}) \geq k - \varepsilon_{i} 
\end{equation}
 for each $i$ in $ \{ 0, 1, \ldots, 4 k - 3 \}$. 
 In a similar fashion, a combined application of \eqref{2kminusdi}
 and Lemma \ref{lemmadistance} gives us that $ d(x_{i}, x_{i+2}) \leq 4 k - d_i - d_{i+1} \leq 4 k - (2 k - f_i) - 
 (2 k - f_{i+1}) = k +\varepsilon_i$. This, together with \eqref{negepsilon}, gives us that 
\begin{equation}\label{boundabs} 
 \left| k - d(x_i, x_{i+2}) \right| \leq \varepsilon_i. 
\end{equation}

 Now, let $\overline{C}$ denote a copy of $C_{2k}$, 
 with its vertex set written so that 
 $V(\overline{C}) = \{ \overline{v}_{0}, \overline{v}_{1}, \ldots, \overline{v}_{2k-1} \}$. 
 We also write $\overline{d}$ to denote the distance function for $\overline{C}$. 
 We then define a quotient mapping 
\begin{equation}\label{rhocolon} 
 \rho\colon V(C_{4k}) \to V(\overline{C})
\end{equation}
 so that 
\begin{equation}\label{rhovalue}
 \rho(v_{j}) = \overline{v}_{j \bmod {2k}}. 
\end{equation} 
 The function defined via \eqref{rhovalue} thus maps vertices in the same antipodal pair of $C_{4k}$
 to the same vertex in the codomain in \eqref{rhocolon}. 
 Now, define
\begin{equation}\label{taucolon}
 \tau\colon V(\overline{C}) \to V(\overline{C})
\end{equation}
 so that 
 $ 
 \tau(\overline{v}_{j}) = \overline{v}_{(j+k) \bmod {2k}} $ 
 for each $j$ in $ \{ 0, 1, \ldots, 2k-1 \}$. 

 The definition in \eqref{rhovalue} gives us that 
\begin{equation}\label{linedrhorho}
 \overline{d}(\rho(u), \rho(v)) = \min\{ d(u, v), 2 k - d(u, v) \} 
\end{equation}
 for all $u, v \in V(C_{4k})$. 
 Now, since $\tau(\rho(v))$ is the antipode of $\rho(v)$ in $C_{2k}$, we find that 
$ \overline{d}(\rho(u), \tau(\rho(v)))
 = k - \overline{d}(\rho(u), \rho(v)) 
 = k - \min\{ d(u, v), 2 k - d(u, v) \}$, i.e., so that 
\begin{equation}\label{abskduv}
 \overline{d}(\rho(u), \tau(\rho(v))) = \left| k - d(u, v) \right|. 
\end{equation}
 Also, as a consequence of \eqref{linedrhorho}, we have that 
\begin{equation}\label{2kminusduv} 
 \overline{d}(\rho(u), \rho(v)) \leq 2 k - d(u, v) 
\end{equation}
 for arbitrary $u, v \in V(C_{4k})$. 

 Now, define
\begin{equation}\label{defineyi}
 y_{i} = \tau^{\left\lfloor \frac{i}{2} \right\rfloor}\left( \rho(x_{i}) \right) 
\end{equation}
 for $i \in \{ 0, 1, \ldots, 4k-1 \}$. 
 Observe that since $\tau$ is an involution, 
 the power of $\tau$ in the right-hand side of \eqref{defineyi} is determined by the residue class of $i$ modulo $4$.
 Also observe that $\tau$ is an isometry, i.e. 
 so that 
 $ \overline{d}\left( \tau(\overline{u}), \tau(\overline{v}) \right) 
 = \overline{d}\left( \overline{u}, \overline{v} \right) $ 
 for all vertices $\overline{u}$ and $\overline{v}$ in $V(\overline{C})$. 
 
 Using the property such that $\tau$ is an isometry together with the relation in \eqref{abskduv}, we find that $\overline{d}(y_i, y_{i + 
 2}) = \overline{d}(\rho(x_i), \tau(\rho(x_{i+2})))  = |k - d(x_{i}, x_{i+2}) |$, so that an application of \eqref{boundabs} 
 then gives us that 
\begin{equation}\label{metricgaptwo}
 \overline{d}(y_{i}, y_{i+2}) \leq \varepsilon_{i}. 
\end{equation}
 From the definition in \eqref{defineyi}, we see that the vanishing of the superscript on the right of \eqref{defineyi} for $i \in \{ 0, 
 1 \}$ gives us that 
\begin{equation}\label{fromyi} 
 \overline{d}(y_0, y_1) = \overline{d}\big( \rho(x_0), \rho(x_1) \big). 
\end{equation}
 A combined application of \eqref{2kminusdi}, \eqref{2kminusduv}, and 
 \eqref{fromyi} allows us to deduce that 
\begin{equation}\label{firstendpoint}
 \overline{d}(y_0, y_1) \leq 2 k - d_0 \leq f_0. 
\end{equation}
Now, since
 $\left\lfloor \frac{4k-2}{2} \right\rfloor = \left\lfloor \frac{4k-1}{2} \right\rfloor = 2 k - 1$, 
 we find that 
\begin{equation}\label{y4kpair}
 y_{4k-2} = \tau^{2k-1}\left( \rho(x_{4k-2}) \right) \quad \text{and} \quad
 y_{4k-1} = \tau^{2k-1}\left( \rho(x_{4k-1}) \right). 
\end{equation}
 Since $\tau$ is an involution, we obtain from \eqref{y4kpair} that 
\begin{equation}\label{y4kagain}
 y_{4k-2} = \tau\left( \rho(x_{4k-2}) \right) \quad \text{and} \quad
 y_{4k-1} = \tau\left( \rho(x_{4k-1}) \right). 
\end{equation}
 Since $\tau$ is an isometry, this and \eqref{y4kagain} together give us that 
\begin{align}
\begin{split}
 \overline{d}(y_{4k-2}, y_{4k-1})
 & = \overline{d}(\tau(\rho(x_{4k-2})), \tau(\rho(x_{4k-1}))) \\
 & = \overline{d}(\rho(x_{4k-2}), \rho(x_{4k-1})). 
\end{split}\label{isometryy4k}
\end{align}
 A combined application of \eqref{2kminusduv} and \eqref{isometryy4k} then gives us that 
\begin{equation}\label{aftertauy4k}
 \overline{d}(\rho(x_{4k-2}), \rho(x_{4k-1})) \leq 2 k - d_{4k-2}. 
\end{equation}
 Now, a combined application of \eqref{2kminusdi}, \eqref{isometryy4k}, and \eqref{aftertauy4k} gives us that 
\begin{equation}\label{secondendpoint} 
 \overline{d}(y_{4k-2}, y_{4k-1}) \leq f_{4k-2}. 
\end{equation}

 Consider the tuple 
\begin{equation}\label{displayW}
 W = \big( y_0, y_2, \ldots, y_{4k-2}, y_{4k-1}, y_{4k-3}, \ldots, y_{1}, y_{0} \big) 
\end{equation}
 consisting of vertices in the underlying set of the metric space $(V(\overline{C}), \overline{d})$, noting the first and final entries of the 
 tuple in \eqref{displayW} are the same. 
 We define the \emph{metric length} $\operatorname{met}(W)$ of $W$ (as opposed to the length of $W$ as a tuple) 
 as the sum of the distances between the consecutive entries of $W$. 
 From \eqref{metricgaptwo}, we find that 
\begin{equation}\label{firstepsilonsum} 
 \sum_{\substack{ 0\leq i \leq 4k-4 \\ \text{$i$ even} }} \overline{d}(y_i, y_{i+2}) 
 \leq \sum_{\substack{0\leq i \leq 4k-4 \\ \text{$i$ even}}} \varepsilon_i 
\end{equation}
 and that 
\begin{equation}\label{secondepsilonsum} 
 \sum_{\substack{ 1 \leq i \leq 4 k - 3 \\ \text{$i$ odd} }} \overline{d}(y_{i+2}, y_{i}) 
 \leq \sum_{\substack{ 1 \leq i \leq 4 k - 3 \\ \text{$i$ odd}}} \varepsilon_{i}. 
\end{equation}
 The definition in \eqref{displayW} and the definition of $\operatorname{met}(W)$ give us that 
\begin{equation}\label{firstmetbound} 
 \operatorname{met}(W) = 
 \overline{d}(y_{4k-2}, y_{4k-1}) 
 + \overline{d}(y_{1}, y_{0}) 
 + \sum_{\substack{ 0\leq i \leq 4k-4 \\ \text{$i$ even} }} \overline{d}(y_i, y_{i+2}) + \sum_{\substack{ 1 \leq i \leq 4 k - 3 \\ \text{$i$ odd} }} \overline{d}(y_{i+2}, y_{i}). 
\end{equation}
 A combined application of \eqref{firstendpoint}, 
 \eqref{secondendpoint}, \eqref{firstepsilonsum}, 
 \eqref{secondepsilonsum}, and \eqref{firstmetbound} gives us that 
\begin{equation}\label{metWabove}
 \operatorname{met}(W) \leq f_0 + f_{4k-2} + \sum_{i=0}^{4k-3} \varepsilon_i.
\end{equation}
 We proceed to determine a lower bound for $\operatorname{met}(W)$. 

 Recall that $\overline{C}$ is an isomorphic copy of $C_{2k}$. Also, recall that $y_{i}$, for a given index $i \in \{ 0, 1, \ldots, 4 k - 1 \}$ and 
 according to the definition in \eqref{defineyi}, is obtained by taking a vertex $x_{i}$ in $V(C_{4k})$ (recalling \eqref{VC4k}), mapping 
 $x_i$ to an element in the codomain $V(\overline{C})$ in \eqref{rhocolon}, and then mapping the resultant vertex to an element in 
 the codomain $V(\overline{C})$ in \eqref{taucolon}, under possibly repeated applications of the involution $\tau$. Also observe that 
\begin{equation}\label{ztauz}
 \tau(z) \neq z
\end{equation}
 for each element $z$	 in the domain $V(\overline{C})$ of $\tau$, i.e., since $\tau$ maps $z$ to the antipodal vertex of $z$. We claim 
 that the sequence $y_{0}$, $y_{1}$, $\ldots$, $y_{4k-1}$ contains at least $k$ distinct vertices of $\overline{C}$, and this is 
 demonstrated below. 

 Fix a vertex $z$ in $V(\overline{C})$. If $y_{i} = z$ for some index $i \in \{ 0, 1, \ldots, 4 k - 1 \}$, then, according to the definition 
 in \eqref{defineyi} (and recalling \eqref{ztauz} and that $\tau$ is an involution), we have that $\rho(x_{i}) \in \{ z, \tau(z) \}$. Each vertex 
 of $\overline{C}$ has precisely two preimages with respect to $\rho$. So, the set $\rho^{-1}(\{ z, \tau(z) \})$ contains precisely four 
 vertices of $C_{4k}$ (recalling the domain in \eqref{rhocolon}). Since the expressions among $x_0$, $x_1$, $\ldots$, and $ x_{4k-1}$ 
 are pairwise unequal (recalling \eqref{VC4k}), any choice of $z \in V(\overline{C})$ can occur at most four times among expressions of 
 the form $y_{i}$ for $i \in \{ 0, 1, \ldots, 4 k - 1 \}$. This shows that 
\begin{equation}\label{atleastky}
 | \{ y_{i} : i \in \{ 0, 1, \ldots, 4 k - 1 \} \} |
 \geq \frac{4k}{4} = k. 
\end{equation}
 From \eqref{displayW} and \eqref{atleastky} together, we find that $W$ contains at 
 least $k$ distinct vertices. 

 We claim that $W$ is of metric length at least $2k-2$. To show this, we begin by forming a closed walk along the edges of $C_{2k}$, by 
 replacing each pair of consecutive entries in $W$ with a shortest path between these vertices, and by forming a walk from the 
 consecutive paths formed. Observe that the length of this resulting walk is equal to $\operatorname{met}(W)$. Now, let $H$ denote 
 the connected subgraph consisting of all vertices and edges traversed by the walk. If $H = C_{2k}$ (letting $\overline{C}$ be identified 
 with $C_{2k}$ for convenience), then each edge of $C_{2k}$ is used at least once, so that the length of the walk is at least $2k$, and 
 hence at least $2k-2$. If $H \neq C_{2k}$, then, since $H$ is (by construction) a connected subgraph of $C_{2k}$, we find that $H$ is a 
 path. Since $W$ contains at least $k$ distinct vertices (recalling \eqref{atleastky}), we then find that $H$ has at least $k - 1$ edges. 
 Each edge of $H$ is a bridge. Since the walk we have constructed is closed and uses every edge of $H$, each edge of $H$ is necessarily
 traversed at least twice. Moreover, since $H$ contains at least $k$ vertices, it has at least $k-1$ edges. Consequently, the length of the 
 walk is at least $2 | E(H) | \geq 2(k-1)$, and thus we thus have the desired relation such that
\begin{equation}\label{metWlower}
 \operatorname{met}(W) \geq 2 k - 2. 
\end{equation} 
 A combined application of \eqref{metWabove} and \eqref{metWlower} then gives us that 
\begin{equation}\label{combinemetW} 
 f_{0} + f_{4k-2} + \sum_{i=0}^{4k-3} \varepsilon_{i} \geq 2 k - 2. 
\end{equation}
 From \eqref{combinemetW} and the definition of $\varepsilon_{i}$ in \eqref{defineepsilon}, we see, 
 with the use of a reindexing argument, that 
\begin{align}
\begin{split}
 \sum_{i=0}^{4k-3} \varepsilon_{i} 
 & = \sum_{i=0}^{4k-3} (f_{i} + f_{i+1} - k) \\ 
 & = 2 \sum_{i=0}^{4k-2} f_{i} - f_{0} - f_{4k-2} - (4k-2) k.
\end{split}\label{reindexing}
\end{align}
 Using \eqref{evaluatespf} and \eqref{reindexing} together, we then obtain that 
\begin{equation}\label{reintroducesp} 
 f_{0} + f_{4k-2} + \sum_{i=0}^{4k-3} \varepsilon_{i} = 2 \operatorname{sp}(f) - (4k-2) k. 
\end{equation}
 The relations in \eqref{combinemetW} and \eqref{reintroducesp} together give us that $ 2 \operatorname{sp}(f) - (4 k - 2) k \geq 2 k - 
 2$, which, in turn, implies that 
\begin{equation}\label{spflower} 
 \operatorname{sp}(f) \geq 2 k^2 - 1. 
\end{equation}
 Since $f$ was chosen as an arbitrary antipodal labeling of $C_{4k}$, we can conclude from \eqref{spflower} that 
 $\operatorname{an}(C_{4k}) \geq 2k^2 - 1$. This and the Juan--Liu upper bound in \eqref{C4kbounds} 
 \cite[Theorem 9]{JuanLiu2012}
 then give us the desired result. 
\end{proof}

\begin{thebibliography}{99}
\bibitem[]{ChartrandErwinZhang2000} Chartrand, Gary; Erwin, David; Zhang, Ping. Radio antipodal colorings of cycles. Congr. Numer. 144 (2000) 129--141
\bibitem[]{ChartrandErwinZhang2002} Chartrand, Gary; Erwin, David; Zhang, Ping. Radio antipodal colorings of graphs. Math. Bohem. 127 (2002) 57--69
\bibitem[]{JuanLiu2012} Juan, Justie Su-Tzu; Liu, Daphne Der-Fen. Antipodal labelings for cycles. Ars Combin. 103 (2012) 81--96
\bibitem[]{KumarPanigrahi2026} Kumar, Kush; Panigrahi, Pratima. Antipodal coloring number of graphs with specified triameter. Hacettepe Journal of Mathematics and Statistics (2026)
\end{thebibliography}
\end{document}
